\documentclass[11pt,a4paper]{ctexart}
\usepackage[color]{tlnote}
\usepackage{needspace}
\title{高数 I（A 类）期中 · 作图七步首遍完整训练卡}
\author{陶理}
\noteheader{先完整做一遍，再谈熟练}
\docdate{2026-11-10}

\begin{document}
\renewcommand{\arraystretch}{1.20}
\noindent{\Large\bfseries 作图七步 · 首遍完整训练卡}\hfill{\small\color{tlblue}11/10 / 11/11 期中倒数第二题}
\par\noindent\rule{\linewidth}{0.4pt}\par\smallskip
\noindent{\small\color{gray}你说「还没完整做过一遍画图、列表不熟」。核查后确认：\textbf{这不是练得少，是练的动作根本不是考的那个动作}。本卡只做一件事——用 50 分钟，让手完整地走完一遍「算 → 填表 → 画图」；\textbf{只新增 1 道题}。}

\section{一、为什么练了 22 道还是不会列表}

\begin{center}\footnotesize
\begin{tabular}{@{}p{3.2cm}p{5.6cm}p{6.2cm}@{}}
\toprule
\textbf{册子} & \textbf{M1 实际在练什么} & \textbf{缺了什么} \\
\midrule
定向加练 M1（14 题） & 斜渐近线、垂直渐近线、驻点、拐点、极值、水平渐近线 & \textbf{无一道要求列表或画图}；方法卡 7 步以「渐近线」收尾，没有「描图」 \\
网络题源 M1（8 题） & 渐近线、极值、拐点、两端渐近线 & 同样全是零件，无完整作图 \\
列表式讲义（5 页） & 6 道例题 + L1/L2 共 \textbf{8 张已复核的五行表} & 是\textbf{示范}，你从没被要求自己填一遍 \\
\bottomrule
\end{tabular}
\end{center}

\noindent{\footnotesize 真题口径（《真题考点索引 2008--2025》）：作图约 \textbf{16/17 套}出现在倒数第二题，10--12 分；卷面常附 $y'$, $y''$，考的是「用导数列表 $\to$ 定性态 $\to$ 画图 $\to$ 找渐近线」的\textbf{流程}，不是硬算导数。笔记明确：\textbf{12 分题必须写全五行，只画图不给表要扣分}。}

\section{二、别背顺序，背「5 个必交件」}

\noindent{\footnotesize 三份材料的七步步序不一致（清单 §7.3 把渐近线放第 6 步、列表式讲义放第 3 步、定向加练 M1 卡则漏掉列表与描图）。\textbf{顺序不影响得分，卷面上有没有这 5 样才影响得分}。计算顺序以明早要默的《考前知识点清单》§7.3 为准，不必纠结。}

\begin{center}\footnotesize
\begin{tabular}{@{}p{0.9cm}p{3.0cm}p{10.7cm}@{}}
\toprule
\textbf{\#} & \textbf{必交件} & \textbf{卷面上必须出现} \\
\midrule
1 & 定义域 & 写成区间，不是只写「$x\ne2$」 \\
2 & 三类渐近线 & 垂直 / 水平 / 斜，\textbf{每类都写「有」或「无」}；两端分别查 \\
3 & \textbf{五行性态表} & $x$ ／ $y'$ ／ 升降 ／ $y''$ ／ 凹凸，\textbf{缺一行扣一行} \\
4 & 关键点坐标 & 极值点、拐点（含「无拐点」结论）、与坐标轴交点 \\
5 & 图 & 先画渐近线（虚线）$\to$ 描点 $\to$ 按表连线；\textbf{曲线必须过已求出的点} \\
\bottomrule
\end{tabular}
\end{center}

\begin{warnbox}[填表六条硬规则（列表式讲义 §1）]
\begin{enumerate}[leftmargin=1.8em,itemsep=1pt]
\item \textbf{固定五行}：$x$（分点与区间交替）／$y'$ 符号／$y$ 升降／$y''$ 符号／凹凸。
\item \textbf{分点四类合并排序}：无定义点、$y'=0$ 驻点、$y'$ 不存在的不可导点、$y''=0$ 或不存在的拐点候选。
\item \textbf{四种记号}：$+$、$-$、$0$、不存在。分点处导数发散写「不存在」，\textbf{不要写 $\infty$}。
\item \textbf{凹凸口径（同济）}：$y''>0$ 记「凹」，$y''<0$ 记「凸」；拐点必须左右变号，$f''=0$ 不变号不算。
\item \textbf{升降行}由 $y'$ 符号读：$+\to\nearrow$，$-\to\searrow$；驻点写「极大 / 极小 + 函数值」，无定义点写「无定义」。
\item \textbf{作图不在表里}：表填完再画；卷面附了 $y'$, $y''$ 时\textbf{不要回头重算导数}，省下的时间给渐近线两端与表。
\end{enumerate}
\end{warnbox}

\noindent{\footnotesize\textbf{可直接抄到考卷的空白模板}（按实际分点数增删列）：}

\begin{center}
\footnotesize\setlength{\tabcolsep}{4pt}
\begin{tabular}{@{}c|ccccccc@{}}
$x$ & $(-\infty,a)$ & $a$ & $(a,b)$ & $b$ & $(b,c)$ & $c$ & $(c,+\infty)$ \\
\hline
$y'$ & & & & & & & \\
$y$ & & & & & & & \\
$y''$ & & & & & & & \\
凹凸 & & & & & & & \\
\end{tabular}
\end{center}

\needspace{10\baselineskip}
\section{三、三遍法：50 分钟走完第一遍（今天）}

\begin{center}\footnotesize
\begin{tabular}{@{}p{1.2cm}p{1.4cm}p{10.2cm}p{1.5cm}@{}}
\toprule
\textbf{遍} & \textbf{时长} & \textbf{只做这一件事} & \textbf{完成} \\
\midrule
1 照抄 & 8 min & 打开《讲义\_函数性态作图\_列表式\_11-09.pdf》§1 空白模板 + §4 例 1（$y=\frac{x^2}{x-1}$）填好的表，\textbf{手抄一遍}，边抄边念五行名字：$x$ / $y'$ / 升降 / $y''$ / 凹凸。目的只有一个：把表的形状刻进手里 & $\square$ \\
2 半独立 & 15 min & 同册 §5 限时卷 \textbf{L1}（$y=\frac{(x-1)^2}{x+1}$）。\textbf{折住标准答案表}，只看它给的「分点 $-3$、$-1$、$1$」提示，自己填五行表 + 画草图，再展开对照。这遍专练「分点 $\to$ 符号 $\to$ 升降/凹凸」 & $\square$ \\
3 全独立 & \textbf{22 min 掐表} & 本卡 §4 的题：从定义域到描图\textbf{一次写完}，五行表 + 图缺一不可。22 min 是《考前知识点清单》§1 给作图题的硬预算，第一次就按考场节奏走 & $\square$ \\
对答案 & 5 min & 折开本卡 §5 标准答案表，逐行对；再勾 §6 的 7 项自检 & $\square$ \\
\bottomrule
\end{tabular}
\end{center}

\begin{keypoint}[为什么是这个顺序]
不熟列表时，\textbf{导数计算会占满工作记忆}，根本顾不上表的结构。所以遍 1 抄形状、遍 2 只练符号映射（分点已给）、遍 3 才放开全流程。三遍都只碰\textbf{同一张表}，不做第二道新题——首遍要的是「完整」，不是「多」。
\end{keypoint}

\vspace{4pt}
\begin{keypoint}[附：三份材料的七步为什么不一样（看一眼就够，不用背）]
\begin{tabular}{@{}p{3.2cm}p{6.8cm}p{4.4cm}@{}}
\textbf{出处} & \textbf{写出来的步序} & \textbf{差在哪} \\
\hline
考前知识点清单 §7.3（明早默它） & 定义域 $\to$ 奇偶 $\to$ 截距 $\to$ $y'$ 表 $\to$ $y''$ 表 $\to$ 渐近线 $\to$ 描图 & \textbf{以此为准}，5 个必交件齐全 \\
作图列表式讲义 §2 & 定义域 $\to$ 奇偶 $\to$ \textbf{渐近线} $\to$ $f'$ $\to$ $f''$ $\to$ \textbf{合并列表} $\to$ 作图 & 渐近线提前只是\textbf{计算顺序}不同；两张表合并成一张五行表 \\
定向加练 M1 方法卡 & 定义域 $\to$ 奇偶 $\to$ 截距 $\to$ $y'$ $\to$ 极值 $\to$ $y''$ $\to$ 渐近线 & \textbf{漏了「列表」和「描图」}，别照它背 \\
\end{tabular}
\par\vspace{3pt}
{\footnotesize 三处并不冲突：一个写的是「先算什么」，一个写的是「最后交什么」。真正固定的是本卡 §2 的 5 个必交件和上面那张五行表模板。}
\end{keypoint}

\newpage
\section{四、独立做（22 min 掐表，先别翻下一页）}

\begin{keypoint}[题目]
求作函数 $y=\dfrac{x^{2}+5}{x-2}$ 的图形（按交大 Q17 形态，\textbf{卷面附导数}，非真题原件）。
\par\smallskip
{\footnotesize 卷面提示：$y'=\dfrac{(x-5)(x+1)}{(x-2)^{2}}$，$y''=\dfrac{18}{(x-2)^{3}}$。\textbf{直接用于列表，不必重算导数}。}
\end{keypoint}

\noindent{\footnotesize\textbf{(1) 定义域 / 奇偶 / 截距：}\blank[11.6cm]}
\par\vspace{10pt}\noindent{\footnotesize\textbf{(2) 三类渐近线（垂直 / 水平 / 斜，两端分别查，无则写「无」）：}}
\par\vspace{4pt}\noindent{\footnotesize\blank[14.4cm]}
\par\vspace{4pt}\noindent{\footnotesize\blank[14.4cm]}

\vspace{8pt}
\noindent{\footnotesize\textbf{(3) 五行性态表}（先写分点，再填五行；分点列按实际个数增删）：}

\begin{center}
\footnotesize\setlength{\tabcolsep}{3pt}
\begin{tabular}{@{}c|ccccccc@{}}
$x$ & \blank[1.5cm] & \blank[1.5cm] & \blank[1.5cm] & \blank[1.5cm] & \blank[1.5cm] & \blank[1.5cm] & \blank[1.5cm] \\
\hline
$y'$ & & & & & & & \\
\hline
$y$ & & & & & & & \\
\hline
$y''$ & & & & & & & \\
\hline
凹凸 & & & & & & & \\
\end{tabular}
\end{center}

\vspace{6pt}
\noindent{\footnotesize\textbf{(4) 关键点坐标（极值点 / 拐点 / 交点）：}\blank[9.6cm]}

\vspace{10pt}
\noindent{\footnotesize\textbf{(5) 作图区}（先画渐近线虚线 $\to$ 描点 $\to$ 按表连线）：}

\begin{center}
\fbox{\parbox[c][6.0cm][t]{14.4cm}{%
  \vspace{3pt}\footnotesize\color{gray} 渐近线用虚线；极值点、截距要落在曲线上；两支曲线不要跨过 $x=2$ 连起来。\par\vfill
  \centering\color{gray}（作图区）\par\vfill
}}
\end{center}

\vspace{6pt}
\noindent{\footnotesize 开始 \blank[1.8cm] \quad 结束 \blank[1.8cm] \quad 实际用时 \blank[1.8cm] \quad 目标 $\le 22$ min}

\newpage
\section{五、标准答案（做完再折开）}

\noindent{\footnotesize 全部数值与符号由 \texttt{Source/\_verify\_graph\_firstpass.py} 用 sympy 独立复核（\textbf{32 项 PASS / 0 FAIL}）。}

\begin{center}\footnotesize
\begin{tabular}{@{}p{2.4cm}p{12.2cm}@{}}
\toprule
\textbf{项} & \textbf{结论} \\
\midrule
定义域 & $(-\infty,2)\cup(2,+\infty)$；非奇非偶 \\
代数变形 & $y=\dfrac{x^{2}+5}{x-2}=x+2+\dfrac{9}{x-2}$（多项式除法直接读出斜渐近线） \\
垂直渐近线 & $x=2$：左极限 $-\infty$、右极限 $+\infty$ \\
水平渐近线 & \textbf{无}（两端 $y\to\pm\infty$） \\
斜渐近线 & $y=x+2$，\textbf{两端都是}这一条 \\
与轴交点 & $(0,-\frac52)$；$x^{2}+5=0$ 无实根 $\Rightarrow$ \textbf{与 $x$ 轴无交点} \\
分点 & $-1$（极大）、$2$（无定义）、$5$（极小），三类合并排序共 3 个 \\
极值 & 极大 $(-1,-2)$；极小 $(5,10)$ \\
拐点 & \textbf{无}。$y''=\frac{18}{(x-2)^{3}}$ 在定义域内无零点；$x=2$ 不在定义域，\textbf{不能标成拐点} \\
\bottomrule
\end{tabular}
\end{center}

\vspace{6pt}
\noindent{\footnotesize\textbf{五行性态表（标准答案）}：}

\begin{center}
\footnotesize\setlength{\tabcolsep}{3pt}
\begin{tabular}{@{}c|ccccccc@{}}
$x$ & $(-\infty,-1)$ & $-1$ & $(-1,2)$ & $2$ & $(2,5)$ & $5$ & $(5,+\infty)$ \\
\hline
$y'$ & $+$ & $0$ & $-$ & 不存在 & $-$ & $0$ & $+$ \\
$y$ & $\nearrow$ & 极大 $-2$ & $\searrow$ & 无定义 & $\searrow$ & 极小 $10$ & $\nearrow$ \\
$y''$ & $-$ & $-$ & $-$ & 不存在 & $+$ & $+$ & $+$ \\
凹凸 & 凸 & 凸 & 凸 & —— & 凹 & 凹 & 凹 \\
\bottomrule
\end{tabular}
\end{center}

\begin{keypoint}[本题两个必踩的给分点]
① $x=2$ 把区间切成两段，两侧 $y''$ 异号但\textbf{都不是拐点}（点不在定义域）；表中该列写「不存在 / 无定义」，不能写 $\infty$。
② 两端斜渐近线是同一条 $y=x+2$，但\textbf{水平渐近线仍要写「无」}——三类逐类交代才算完整。
\end{keypoint}

\section{六、7 项自检（交卷前 60 秒）}

\begin{enumerate}[leftmargin=1.9em,itemsep=1pt]
\item[$\square$] 定义域写成\textbf{区间}，不只是「$x\ne2$」。
\item[$\square$] 垂直渐近线两侧极限都写了；只一侧发散也算。
\item[$\square$] 水平渐近线\textbf{两端各查一次}，本题写「无」。
\item[$\square$] 斜渐近线\textbf{两端各查一次}，写出 $k$ 与 $b$ 的求法或多项式部分。
\item[$\square$] 五行表齐：$x$ / $y'$ / 升降 / $y''$ / 凹凸，分点四类合并排序。
\item[$\square$] 无定义点列写「不存在 / 无定义」，\textbf{没把 $x=2$ 标成拐点}。
\item[$\square$] 图上渐近线是虚线，曲线\textbf{经过} $(-1,-2)$、$(5,10)$、$(0,-\frac52)$，且两支没跨 $x=2$ 相连。
\end{enumerate}

\vspace{6pt}
\noindent{\footnotesize\color{gray} 材料入口：\textbf{本卡}（首遍完整训练）$\mid$ 《讲义\_函数性态作图\_列表式\_11-09.pdf》5 页（8 张已复核的五行表，遍 1 / 遍 2 用）$\mid$ 《考前知识点清单》§7.3（七步与 22 min 预算）。本卡独立 PDF，不改任何既有包绝对页码；不新增题册，全卡只有 1 道新题。}

\end{document}
